\begin{lemma}[Agreement of the kernel and spectral realizations on a smooth core]
Fix $n \geq 1$ and $j \in \{1,\ldots,n\}$.  Let $f : \mathbb{R}^n \to \mathbb{R}$
be smooth with compact support, and let $x \in \mathbb{R}^n$.  For $s > 0$ write
\[
  r = e^{-s},\qquad \tau = \sqrt{1-e^{-2s}},\qquad
  (A_sg)(x) = \int_{\mathbb{R}^n} g(rx+\tau z)\,d\gamma_n(z)
\]
for the Mehler average of \noderef{MehlerFormula}, let $D_j = 2^{-1/2}\partial_j$
be the energy-normalized partial derivative of \noderef{GaussD}, and put
\[
  \Phi_j(x) = \frac{1}{\sqrt\pi}\int_0^\infty
      e^{-s}\,\bigl(A_s(D_jf)\bigr)(x)\,\frac{ds}{\sqrt s}
    = \frac{1}{\sqrt\pi}\int_0^\infty
      e^{-s}\,\frac{1}{\sqrt2}\bigl(A_s(\partial_jf)\bigr)(x)\,\frac{ds}{\sqrt s} ,
\]
the function of \noderef{RieszSemigroupPointwise}.  Let $K_{\gamma,j}$ be the
normalized Gaussian Riesz kernel of \noderef{GaussRieszKernel}.  Then:
\begin{enumerate}
\item for every $\varepsilon > 0$ the function $y \mapsto K_{\gamma,j}(x,y)f(y)$
is Lebesgue integrable on the exterior region $\{y : |x-y| > \varepsilon\}$; and
\item
\[
  \lim_{\varepsilon\downarrow0}
    \int_{|x-y|>\varepsilon}K_{\gamma,j}(x,y)\,f(y)\,dy
  \;=\; \Phi_j(x) ,
\]
the limit being taken along $\varepsilon \to 0$ with $\varepsilon > 0$.
\end{enumerate}

This is the reference paper's appendix lemma: the truncated kernel integrals of
a smooth compactly supported input converge, at \emph{every} $x$ and not merely
almost everywhere, to the Gaussian Riesz transform of that input, the
right-hand side being read as the continuous semigroup representative.  That
reading is exactly what $\Phi_j$ is: by \noderef{RieszSemigroupPointwise} the
integral defining $\Phi_j(x)$ converges absolutely at every $x$ (clause (1)),
$\Phi_j$ is continuous on $\mathbb{R}^n$ (clause (2)), and $\Phi_j$ represents
the $j$th coordinate of $R_\gamma u$, where $u \in L^2(\gamma_n)$ is the class
of $f$ (clause (3)).  The every-$x$ strength of the present statement is what
later nodes consume, so it is not weakened to an almost-everywhere assertion;
it is available because both sides are honest pointwise quantities, the
left-hand side by clause (1) above and the right-hand side by the continuity of
$\Phi_j$.

Every constant appearing in the proof may depend on $n$, on $j$, on $x$ and on
$f$.  As the appendix header instructs, these constants serve only to identify
two realizations of one operator, and nothing below is uniform in $n$ or feeds
any endpoint constant.
\end{lemma}
\begin{proof}
Throughout, $\lambda$ is Lebesgue measure on $\mathbb{R}^n$, $\gamma_n$ is the
Gaussian measure of \noderef{GaussMeasure}, and $n \geq 1$, $j$, $f$ and $x$ are
fixed.  For $s > 0$ we keep the abbreviations
\[
  r = e^{-s} \in (0,1),\quad \sigma = 1-r^2 = 1-e^{-2s} \in (0,1),
  \quad \tau = \sqrt\sigma, \quad d = (1-r)x ,
\]
and we write
\[
  p_s(x,y) = \frac{1}{\pi^{n/2}\sigma^{n/2}}
      \exp\Bigl(-\frac{|rx-y|^2}{\sigma}\Bigr),
  \qquad
  q_\sigma(z) = (\pi\sigma)^{-n/2}e^{-|z|^2/\sigma},
  \qquad
  V_{j,\sigma}(z) = \frac{z_j}{\sigma}q_\sigma(z) ,
\]
so that, exactly as in \noderef{PVKernelConstantCancellation}, with $z = y-x$ we
have $z+d = y-rx$, $p_s(x,y) = q_\sigma(z+d)$ and
\begin{equation*}
  D_j^xp_s(x,y) := \frac{1}{\sqrt2}\,\partial_{x_j}p_s(x,y)
    = \sqrt2\,r\,\frac{y_j-rx_j}{\sigma}\,p_s(x,y)
    = \sqrt2\,r\,V_{j,\sigma}(z+d) , \tag{0}
\end{equation*}
the middle equality being \noderef{MehlerKernelXDerivative}.  Finally put
\[
  \Phi^{(s)}(w) = \int_{\mathbb{R}^n}p_s(w,y)f(y)\,dy, \qquad w \in \mathbb{R}^n .
\]

\emph{Step 0: elementary bounds on $f$.}
Let $S = \operatorname{supp}f$, a compact set, and let
$M_0 = \sup_{\mathbb{R}^n}|f|$ and $L_f = \sup_{\mathbb{R}^n}|\nabla f|$.  Both
are finite: $f$ and $|\nabla f|$ are continuous and vanish off $S$ (at a point
outside $S$ the function $f$ vanishes on a neighbourhood, hence so do all its
derivatives), so the extreme value theorem on $S$ (shared context) bounds them.
Also $\lambda(S) < \infty$, a compact set having finite Lebesgue measure (shared
context).  By the mean value inequality applied along the segment from $x$ to
$y$ --- $t \mapsto f(x+t(y-x))$ is smooth with derivative
$\nabla f(x+t(y-x))\cdot(y-x)$, so
$|f(y)-f(x)| \leq \int_0^1|\nabla f(x+t(y-x))|\,|y-x|\,dt$ --- we get
\begin{equation*}
  |f(y)-f(x)| \leq L_f\,|y-x| \qquad \text{for all } y \in \mathbb{R}^n . \tag{1}
\end{equation*}
Finally $f$ is bounded and continuous and $\gamma_n$ is a probability measure
(\noderef{GaussIsProb}), so $f$ is square $\gamma_n$-integrable and determines a
class $u \in L^2(\gamma_n)$; this is the $u$ of the statement, and it is the
hypothesis under which \noderef{RieszSemigroupPointwise} is applied below.

\emph{Step 1: $\gamma_n$ and $\lambda$ have the same null sets, and $\gamma_n$
has full support.}
By \noderef{GaussMeasure}, $\gamma_n$ is $\lambda$ weighted by the density
$\rho(y) = \pi^{-n/2}e^{-|y|^2}$, i.e.\ $\gamma_n(A) = \int_A\rho\,d\lambda$ for
measurable $A$.  The density is finite and strictly positive at every point.

If $\lambda(A) = 0$ then $\int_A\rho\,d\lambda = 0$, so $\gamma_n(A) = 0$.
Conversely suppose $\gamma_n(A) = 0$ and put $A_m = A\cap\{\rho>1/m\}$ for
$m \in \mathbb{N}$.  Then
$\lambda(A_m)/m \leq \int_{A_m}\rho\,d\lambda \leq \gamma_n(A) = 0$, so
$\lambda(A_m) = 0$; since $\rho > 0$ everywhere we have $A = \bigcup_m A_m$, and
countable subadditivity (shared context) gives $\lambda(A) = 0$.  So the two
measures have the same null sets.

If $U \subseteq \mathbb{R}^n$ is open and nonempty it contains a ball, so
$\lambda(U) > 0$, hence $\gamma_n(U) > 0$ by the previous paragraph.
Consequently:
\begin{equation*}
  \parbox{0.86\textwidth}{if $g_1,g_2 : \mathbb{R}^n \to \mathbb{R}$ are
  continuous and $g_1 = g_2$ $\gamma_n$-almost everywhere, then $g_1 = g_2$
  everywhere.} \tag{2}
\end{equation*}
Indeed, if $g_1(w_0) \neq g_2(w_0)$ then $|g_1-g_2| > 0$ on some open ball $U$
around $w_0$ by continuity of $g_1-g_2$, and $\gamma_n(U) > 0$ contradicts
$\gamma_n(\{g_1 \neq g_2\}) = 0$.

\emph{Step 2: the kernel form of the semigroup derivative.}
Fix $s > 0$.

\emph{(a) $\Phi^{(s)}$ is well defined and continuous on $\mathbb{R}^n$.}
For each $w$ the function $y \mapsto p_s(w,y)f(y)$ is continuous, vanishes off
$S$ and is bounded in absolute value by $\pi^{-n/2}\sigma^{-n/2}M_0$, because
the exponential factor of $p_s$ is at most $1$; since $\lambda(S) < \infty$ it
is Lebesgue integrable, so $\Phi^{(s)}(w)$ is a well-defined real number.  If
$w_k \to w$ then $p_s(w_k,y)f(y) \to p_s(w,y)f(y)$ for every $y$, by continuity
of $p_s$ in its first argument, and all these functions are dominated by the
integrable function $\pi^{-n/2}\sigma^{-n/2}M_0\mathbf 1_S$; dominated
convergence (shared context) gives $\Phi^{(s)}(w_k) \to \Phi^{(s)}(w)$.  As
$\mathbb{R}^n$ is a metric space, $\Phi^{(s)}$ is continuous.

\emph{(b) $\Phi^{(s)} = A_sf$ everywhere on $\mathbb{R}^n$.}
Let $\tilde u$ denote the function underlying the class $u$ of Step 0, the same
representative that \noderef{MehlerFormula} and \noderef{MehlerLebesgueKernel}
use.  By clause (2) of \noderef{MehlerFormula} applied to $s$ and $u$, and
clause (2) of \noderef{MehlerLebesgueKernel} applied to the same $s$ and $u$,
there is a $\gamma_n$-null set off which both
\[
  (P_su)(w) = \int_{\mathbb{R}^n}\tilde u(rw+\tau z)\,d\gamma_n(z)
  \qquad\text{and}\qquad
  (P_su)(w) = \int_{\mathbb{R}^n}p_s(w,y)\,\tilde u(y)\,dy
\]
hold; hence for $\gamma_n$-almost every $w$,
\begin{equation*}
  \int_{\mathbb{R}^n}p_s(w,y)\,\tilde u(y)\,dy
    = \int_{\mathbb{R}^n}\tilde u(rw+\tau z)\,d\gamma_n(z) . \tag{3}
\end{equation*}
Now replace $\tilde u$ by $f$ on both sides of $(3)$, for \emph{every} $w$.  Let
$N = \{\tilde u \neq f\}$; then $\gamma_n(N) = 0$, since $\tilde u$ is a
representative of the class of $f$, and therefore $\lambda(N) = 0$ by Step 1.
On the left, the integrands $p_s(w,\cdot)\tilde u$ and $p_s(w,\cdot)f$ differ
only on $N$, a $\lambda$-null set, so the two Lebesgue integrals agree for every
$w$ (and in particular the second is absolutely convergent by (a)).  On the
right, fix $w$ and note that
$\{z : rw+\tau z \in N\} = \tau^{-1}(N - rw)$, which is $\lambda$-null because
$\lambda$ is translation invariant (shared context) and the dilation
$v \mapsto \tau^{-1}v$ multiplies $\lambda$ by $\tau^{-n}$; hence it is
$\gamma_n$-null by Step 1, and the two $\gamma_n$-integrals agree for every $w$
(the one with $f$ being absolutely convergent by clause (1) of
\noderef{SemigroupGradientCommute}).  Therefore
\[
  \Phi^{(s)}(w) = (A_sf)(w) \qquad\text{for } \gamma_n\text{-almost every } w .
\]
Both sides are continuous: the left by (a), and the right because clause (4) of
\noderef{SemigroupGradientCommute} asserts that $A_sf$ is differentiable at
every point, hence continuous there.  By $(2)$ of Step 1 they agree at every
$w$.

\emph{(c) the kernel form.}  Set
$G(w,y) = 2r(y_j-rw_j)\sigma^{-1}p_s(w,y)$; by
\noderef{MehlerKernelXDerivative} the map $w \mapsto p_s(w,y)$ is differentiable
at every $w$ and $2^{-1/2}G(w,y) = D_j^xp_s(w,y)$, i.e.\
$G(w,y) = \partial_{w_j}p_s(w,y)$ --- multiply that node's displayed identity by
$\sqrt2$ and use $\sqrt2\cdot\sqrt2 = 2$.  The function $G$ is continuous on
$\mathbb{R}^n\times\mathbb{R}^n$, being a product and composition of continuous
functions.

We claim
\begin{equation*}
  \lim_{t\to0}\frac{\Phi^{(s)}(x+te_j)-\Phi^{(s)}(x)}{t}
    = \int_{\mathbb{R}^n}G(x,y)f(y)\,dy , \tag{4}
\end{equation*}
$e_j$ being the $j$th standard basis vector.  Let $\eta > 0$.  The set
$\overline B(x,1)\times S$ is compact, so $G$ is uniformly continuous on it
(Heine--Cantor); choose $\delta \in (0,1]$ with
$|G(w,y)-G(x,y)| \leq \eta$ whenever $w \in \overline B(x,1)$,
$|w-x| \leq \delta$ and $y \in S$.  Fix $y \in S$ and $t$ with
$0 < |t| \leq \delta$.  The function $t' \mapsto p_s(x+t'e_j,y)$ is
differentiable on $\mathbb{R}$ with derivative $t' \mapsto G(x+t'e_j,y)$, which
is continuous; so by the fundamental theorem of calculus
\[
  p_s(x+te_j,y) - p_s(x,y) = \int_0^t G(x+t'e_j,y)\,dt' ,
\]
and subtracting $t\,G(x,y) = \int_0^tG(x,y)\,dt'$,
\[
  \bigl|p_s(x+te_j,y)-p_s(x,y)-t\,G(x,y)\bigr| \leq |t|\,\eta ,
\]
because every $x+t'e_j$ with $t'$ between $0$ and $t$ lies in
$\overline B(x,1)$ at distance at most $\delta$ from $x$.  Multiplying by
$f(y)$, which is bounded by $M_0$ and vanishes off $S$, and integrating over
$y$ --- the three functions $y \mapsto p_s(x+te_j,y)f(y)$,
$y \mapsto p_s(x,y)f(y)$ and $y \mapsto G(x,y)f(y)$ are continuous, vanish off
$S$ and are bounded, hence integrable, the last because $G(x,\cdot)$ is
continuous --- we obtain
\[
  \Bigl|\Phi^{(s)}(x+te_j)-\Phi^{(s)}(x)-t\int G(x,y)f(y)\,dy\Bigr|
    \leq |t|\,\eta\,M_0\,\lambda(S) .
\]
Dividing by $|t|$ shows that the difference quotient in $(4)$ differs from
$\int G(x,y)f(y)\,dy$ by at most $\eta M_0\lambda(S)$ for every $t$ with
$0 < |t| \leq \delta$.  As $\eta > 0$ was arbitrary, $(4)$ follows.

On the other hand, by (b) the functions $\Phi^{(s)}$ and $A_sf$ coincide, so the
limit in $(4)$ is also the derivative at $t=0$ of $t \mapsto (A_sf)(x+te_j)$.
By clause (4) of \noderef{SemigroupGradientCommute}, $A_sf$ is differentiable at
$x$ and the value of its derivative at $e_j$ --- which is precisely that
derivative along the line --- equals $e^{-s}\int(\partial_jf)(rx+\tau z)\,
d\gamma_n(z) = e^{-s}\bigl(A_s(\partial_jf)\bigr)(x)$.  Hence
\[
  \int_{\mathbb{R}^n}\partial_{x_j}p_s(x,y)\,f(y)\,dy
    = e^{-s}\bigl(A_s(\partial_jf)\bigr)(x) .
\]
Multiplying by $2^{-1/2}$ and using $D_j^xp_s = 2^{-1/2}\partial_{x_j}p_s$ and
$D_j = 2^{-1/2}\partial_j$ (\noderef{GaussD}),
\begin{equation*}
  \int_{\mathbb{R}^n}D_j^xp_s(x,y)\,f(y)\,dy
    = e^{-s}\,\frac{1}{\sqrt2}\bigl(A_s(\partial_jf)\bigr)(x)
    = e^{-s}\bigl(A_s(D_jf)\bigr)(x)
    =: \Theta(s) . \tag{5}
\end{equation*}
The integral on the left is absolutely convergent, since
$|D_j^xp_s(x,\cdot)|$ is Lebesgue integrable by clause (1) of
\noderef{PVKernelConstantCancellation} and $|f| \leq M_0$.

\emph{Step 3: the subtracted bound.}
We claim that there is a constant $C_n$, depending only on $n$, with
\begin{equation*}
  \int_{\mathbb{R}^n}\bigl|D_j^xp_s(x,y)\bigr|\,|f(y)-f(x)|\,dy
    \leq C_n\,L_f\,e^{-s}\,(1+|x|)
  \qquad\text{for every } s > 0 . \tag{6}
\end{equation*}
First record two Gaussian moment identities.  The dilation $v \mapsto cv$
multiplies $\lambda$ by $c^n$, so for every Borel
$h : \mathbb{R}^n \to [0,+\infty]$ and every $c > 0$,
$\int h(z)\,dz = c^n\int h(cv)\,dv$; applying this with $c = \sqrt\sigma$ and
$h = |z_j|q_\sigma$, respectively $h = |z_j|\,|z|\,q_\sigma$, and using
$q_\sigma(\sqrt\sigma v) = (\pi\sigma)^{-n/2}e^{-|v|^2}$ together with
\noderef{GaussMeasure},
\begin{equation*}
  \int_{\mathbb{R}^n}|z_j|\,q_\sigma(z)\,dz = \sqrt\sigma\,a_n ,
  \qquad
  \int_{\mathbb{R}^n}|z_j|\,|z|\,q_\sigma(z)\,dz = \sigma\,b_n , \tag{7}
\end{equation*}
where $a_n = \int|v_j|\,d\gamma_n(v)$ and $b_n = \int|v_j|\,|v|\,d\gamma_n(v)$
are finite, since $|v_j| \leq |v| \leq 1+|v|^2$ and $|v_j||v| \leq |v|^2$ are
dominated by polynomials and the Gaussian exponential moments item of the
shared context (with $t = 0$) applies.

Now use $(1)$ and $(0)$, substitute $y = x+z$ and then $w = z+d$ (two
translations, which preserve $\lambda$; shared context), and use
$|z| = |w-d| \leq |w|+|d|$:
\begin{align*}
  \int_{\mathbb{R}^n}|D_j^xp_s(x,y)|\,|f(y)-f(x)|\,dy
    &\leq L_f\sqrt2\,r\int_{\mathbb{R}^n}|V_{j,\sigma}(z+d)|\,|z|\,dz\\
    &\leq \frac{\sqrt2\,r\,L_f}{\sigma}
        \int_{\mathbb{R}^n}|w_j|\,\bigl(|w|+|d|\bigr)\,q_\sigma(w)\,dw\\
    &= \frac{\sqrt2\,r\,L_f}{\sigma}\bigl(\sigma b_n + |d|\sqrt\sigma a_n\bigr)
     = \sqrt2\,L_f\,r\Bigl(b_n + a_n\frac{|d|}{\sqrt\sigma}\Bigr) ,
\end{align*}
by $(7)$.  Since $|d| = (1-r)|x|$ and $\sigma = (1-r)(1+r)$ with $0 < r < 1$,
\[
  \frac{|d|}{\sqrt\sigma} = |x|\sqrt{\frac{1-r}{1+r}} \leq |x| ,
\]
so the right-hand side is at most
$\sqrt2\,L_f\,e^{-s}(b_n+a_n|x|) \leq C_nL_fe^{-s}(1+|x|)$ with
$C_n = \sqrt2\max(a_n,b_n)$.  This is $(6)$.

\emph{Step 4: Fubini in $(y,s)$, and the subtracted identity.}
We first record the time-integral representation of the kernel: for every
$y \in \mathbb{R}^n$,
\begin{equation*}
  K_{\gamma,j}(x,y)
    = \frac{1}{\sqrt\pi}\int_0^\infty \frac{D_j^xp_s(x,y)}{\sqrt s}\,ds ,
  \tag{8}
\end{equation*}
with the convention --- the one \noderef{GaussRieszKernel} uses for its own
defining integral --- that an integral of a non-integrable integrand is $0$, and
with the $s$-integrand Lebesgue integrable on $(0,\infty)$ exactly when the
integrand of that defining integral is Lebesgue integrable on $(0,1)$.  Indeed,
by $(0)$ and $p_s(x,y) = \pi^{-n/2}(1-r^2)^{-n/2}e^{-|rx-y|^2/(1-r^2)}$,
\[
  \frac{1}{\sqrt\pi}\,\frac{D_j^xp_s(x,y)}{\sqrt s}
    = \frac{\sqrt2}{\pi^{(n+1)/2}}\,
      \frac{r\,(y_j-rx_j)}{\sqrt s\,(1-r^2)^{n/2+1}}
      \exp\Bigl(-\frac{|rx-y|^2}{1-r^2}\Bigr) ,
\]
while \noderef{GaussRieszKernel} gives
$K_{\gamma,j}(x,y) = \sqrt2\,\pi^{-(n+1)/2}\int_{(0,1)}h(\varrho)\,d\varrho$
with
\[
  h(\varrho) = \frac{y_j-\varrho x_j}{(1-\varrho^2)^{n/2+1}\sqrt{-\log\varrho}}
      \exp\Bigl(-\frac{|\varrho x-y|^2}{1-\varrho^2}\Bigr) .
\]
The map $\psi(s) = e^{-s}$ is a continuously differentiable bijection of
$(0,\infty)$ onto $(0,1)$ with $|\psi'(s)| = e^{-s} \neq 0$ and inverse
$\varrho\mapsto-\log\varrho$, so that $\sqrt{-\log\psi(s)} = \sqrt s$.  The
one-dimensional change-of-variables formula for such a bijection, which inserts
the factor $|\psi'|$, applied first to the nonnegative Borel function $|h|$ (an
identity in $[0,+\infty]$, so that integrability of one side is equivalent to
integrability of the other) and then to $h$ itself, gives
$\int_{(0,1)}h(\varrho)\,d\varrho = \int_{(0,\infty)}h(e^{-s})e^{-s}\,ds$; and
\[
  h(e^{-s})\,e^{-s}
    = \frac{r\,(y_j-rx_j)}{(1-r^2)^{n/2+1}\sqrt s}
      \exp\Bigl(-\frac{|rx-y|^2}{1-r^2}\Bigr) ,
\]
which is $\pi^{(n+1)/2}2^{-1/2}$ times the displayed integrand above.
Multiplying by $\sqrt2\pi^{-(n+1)/2}$ gives $(8)$; and when neither integrand is
integrable both sides of $(8)$ are $0$ by the convention.

Now put $F(y,s) = s^{-1/2}D_j^xp_s(x,y)\bigl(f(y)-f(x)\bigr)$ on
$\mathbb{R}^n\times(0,\infty)$.  It is continuous in $(y,s)$, hence
product-measurable.  By Tonelli (shared context), $(6)$ and
\noderef{SubordinationWeightIntegral},
\[
  \iint |F|
    \leq \int_0^\infty \frac{C_nL_f(1+|x|)e^{-s}}{\sqrt s}\,ds
    = C_nL_f(1+|x|)\sqrt\pi < \infty ,
\]
so $F$ is integrable on the product and Fubini (shared context) applies:
\begin{equation*}
  \int_{\mathbb{R}^n}\Bigl(\int_0^\infty F(y,s)\,ds\Bigr)dy
    = \int_0^\infty\Bigl(\int_{\mathbb{R}^n}F(y,s)\,dy\Bigr)ds , \tag{9}
\end{equation*}
the inner integral on the left being Lebesgue integrable in $y$ and the inner
integral on the right Lebesgue integrable in $s$.

The inner integral on the left is, by $(8)$,
$(f(y)-f(x))\sqrt\pi\,K_{\gamma,j}(x,y)$ for every $y$ --- at the $y$ for which
the $s$-integrand is integrable this is $(8)$, and at the remaining $y$ both
sides are $0$ by the convention.  So
$y \mapsto K_{\gamma,j}(x,y)(f(y)-f(x))$ is Lebesgue integrable on
$\mathbb{R}^n$.  The inner integral on the right is, for each $s > 0$,
\[
  \int_{\mathbb{R}^n}F(y,s)\,dy
    = \frac{1}{\sqrt s}\Bigl(\int D_j^xp_s(x,y)f(y)\,dy
        - f(x)\int D_j^xp_s(x,y)\,dy\Bigr)
    = \frac{\Theta(s)}{\sqrt s} ,
\]
splitting the integral --- legitimate because both pieces are absolutely
convergent, by clause (1) of \noderef{PVKernelConstantCancellation} and
$|f| \leq M_0$ --- and then using $(5)$ for the first piece and clause (1) of
\noderef{PVKernelConstantCancellation} again, which says
$\int D_j^xp_s(x,y)\,dy = 0$, for the second.  Therefore $(9)$ reads
\[
  \sqrt\pi\int_{\mathbb{R}^n}K_{\gamma,j}(x,y)\bigl(f(y)-f(x)\bigr)\,dy
    = \int_0^\infty \frac{\Theta(s)}{\sqrt s}\,ds
    = \sqrt\pi\,\Phi_j(x) ,
\]
the last equality being the definition of $\Phi_j$ together with $(5)$; the
$s$-integral converges absolutely, which is also clause (1) of
\noderef{RieszSemigroupPointwise}.  Dividing by $\sqrt\pi$,
\begin{equation*}
  \int_{\mathbb{R}^n}K_{\gamma,j}(x,y)\bigl(f(y)-f(x)\bigr)\,dy = \Phi_j(x) ,
  \tag{10}
\end{equation*}
the left-hand side being an absolutely convergent Lebesgue integral.  This is
the appendix's subtracted identity.

\emph{Step 5: splitting the truncated integrals, and the limit.}
Fix $\varepsilon > 0$ and write
$\Omega_\varepsilon = \{y : |x-y| > \varepsilon\}$.  By clause (2) of
\noderef{PVKernelConstantCancellation}, $K_{\gamma,j}(x,\cdot)$ is Lebesgue
integrable on $\Omega_\varepsilon$; since $f$ is continuous with
$|f| \leq M_0$, the product $K_{\gamma,j}(x,\cdot)f$ is measurable and dominated
by $M_0|K_{\gamma,j}(x,\cdot)|$ there, hence integrable on $\Omega_\varepsilon$.
This is clause (1) of the lemma.  By $(10)$ the function
$K_{\gamma,j}(x,\cdot)(f-f(x))$ is integrable on $\mathbb{R}^n$, hence on
$\Omega_\varepsilon$, and therefore
\begin{equation*}
  \int_{\Omega_\varepsilon}K_{\gamma,j}(x,y)f(y)\,dy
    = \int_{\Omega_\varepsilon}K_{\gamma,j}(x,y)\bigl(f(y)-f(x)\bigr)\,dy
      + f(x)\int_{\Omega_\varepsilon}K_{\gamma,j}(x,y)\,dy , \tag{11}
\end{equation*}
all three integrals being absolutely convergent.

Let $\varepsilon_k\downarrow0$ be any sequence of positive reals.  The second
term of $(11)$ tends to $0$ by clause (3) of
\noderef{PVKernelConstantCancellation}.  For the first, the functions
$\mathbf 1_{\Omega_{\varepsilon_k}}K_{\gamma,j}(x,\cdot)(f-f(x))$ converge
pointwise on $\mathbb{R}^n\setminus\{x\}$ --- hence $\lambda$-almost everywhere,
a single point being $\lambda$-null --- to
$K_{\gamma,j}(x,\cdot)(f-f(x))$, and they are dominated in absolute value by
that same integrable function; dominated convergence (shared context) gives
\[
  \int_{\Omega_{\varepsilon_k}}K_{\gamma,j}(x,y)(f(y)-f(x))\,dy
    \longrightarrow
  \int_{\mathbb{R}^n}K_{\gamma,j}(x,y)(f(y)-f(x))\,dy = \Phi_j(x)
\]
by $(10)$.  Hence
$\int_{\Omega_{\varepsilon_k}}K_{\gamma,j}(x,y)f(y)\,dy \to \Phi_j(x)$.  Since
the sequence $(\varepsilon_k)$ was arbitrary and $\mathbb{R}$ is a metric space,
so that convergence along $\varepsilon\downarrow0$ is equivalent to convergence
along every such sequence, clause (2) follows.
\end{proof}
